\chapter{NSGA-II : principe et implémentation}\label{ch:nsga}
\iterbanner{V1}

\section{Pourquoi un algorithme génétique}
La note d'intention retient NSGA-II (Deb et al., 2002) pour trois raisons : il ne demande aucun gradient, il est robuste aux discontinuités (en SER, une facette qui passe de l'ombre à la lumière fait sauter l'objectif), et il produit un \emph{ensemble} de solutions à chaque génération, donc une approximation de tout le front en une seule exécution. C'est la référence classique de l'optimisation multi-objectif, disponible dans pymoo.

\section{Le principe}
Un algorithme génétique fait évoluer une \textbf{population} de $N$ solutions. À chaque génération :
\begin{enumerate}
\item on \textbf{sélectionne} des parents, en favorisant les meilleurs ;
\item on les \textbf{croise} et on les \textbf{mute} pour créer $N$ enfants ;
\item on garde les $N$ meilleurs parmi parents et enfants : c'est la \textbf{survie élitiste}.
\end{enumerate}
Tout le problème, en multi-objectif, est de définir « meilleur » quand il n'y a pas d'ordre total. NSGA-II répond par deux critères, dans cet ordre : le \textbf{rang de non-dominance}, puis la \textbf{distance d'encombrement}.

\begin{figure}[htbp]
\centering
\begin{tikzpicture}[node distance=0.45cm and 0.6cm]
\node[flow] (p) {population $P_t$\\ ($N$ profils)};
\node[flow,right=of p] (s) {sélection\\ par tournoi};
\node[flow,right=of s] (v) {croisement SBX\\ mutation polynomiale\\ arrondi, doublons};
\node[flow,right=of v] (e) {$N$ enfants\\ \textbf{simulés}};
\node[flow,below=1.2cm of v,text width=6cm] (r) {$P_t\cup Q_t$ ($2N$) : tri non dominé contraint,\\ remplissage front par front, encombrement};
\node[flow,left=of r] (n) {$P_{t+1}$};
\draw[arr] (p) -- (s); \draw[arr] (s) -- (v); \draw[arr] (v) -- (e);
\draw[arr] (e) |- (r); \draw[arr] (r) -- (n); \draw[arr] (n) -| (p);
\end{tikzpicture}
\caption{Une génération de NSGA-II. Seuls les enfants coûtent des simulations ; les doublons sont éliminés avant d'être simulés.}\label{fig:nsga-cycle}
\end{figure}

\section{Les deux critères de classement}
\paragraph{Le tri non dominé.} On répartit la population en fronts successifs $\mathcal F_1,\mathcal F_2,\dots$ : $\mathcal F_1$ contient les solutions non dominées, $\mathcal F_2$ celles qui ne sont dominées que par $\mathcal F_1$, etc. Le rang d'une solution est l'indice de son front. L'algorithme de Deb compte, pour chaque solution, combien d'autres la dominent, puis « épluche » les fronts : $O(MN^2)$ comparaisons.

\paragraph{La distance d'encombrement.} À rang égal, on préfère les solutions isolées, pour bien répartir la population le long du front. Pour chaque objectif, on trie le front, et on ajoute à chaque point la distance normalisée entre ses deux voisins ; les extrémités reçoivent $+\infty$ pour être toujours conservées :
\[
d_i=\sum_{k=1}^{M}\frac{f_k^{(i+1)}-f_k^{(i-1)}}{f_k^{\max}-f_k^{\min}}.
\]

\paragraph{Les contraintes : règle de Deb.} Pour comparer deux solutions : une faisable bat toujours une infaisable ; entre deux infaisables, la plus petite violation $\mathrm{CV}$ gagne ; entre deux faisables, on applique la dominance. Il n'y a ni pénalité ni coefficient à régler.

\begin{codebox}[ : \code{optim/nsga2.py}, classement contraint]
\begin{lstlisting}
def constrained_ranks(F, CV):
    feas = np.where(CV <= 0)[0]
    r = 0
    for front in non_dominated_sort(F[feas]):        # faisables : fronts de Pareto
        ids = feas[front]
        rank[ids] = r
        crowd[ids] = crowding_distance(F[ids])
        r += 1
    infeas = np.where(CV > 0)[0]
    for i in infeas[np.argsort(CV[infeas])]:         # infaisables : apres, par violation croissante
        rank[i] = r
        r += 1
    return rank, crowd
\end{lstlisting}
\end{codebox}

\section{Les opérateurs, sur des variables entières}
Nos variables sont des indices de 0 à 7. Comme pymoo, on les relâche en réels dans $[-0{,}5;\ 7{,}5]$, on applique des opérateurs continus, puis on arrondit.

\paragraph{Croisement SBX} (\emph{simulated binary crossover}). Pour chaque variable, avec probabilité $\frac12$, les deux enfants sont placés symétriquement autour de la moyenne des parents, à une distance $\beta$ fois leur écart, où $\beta$ suit une loi qui se concentre près de 1 quand $\eta_c$ est grand :
\[
\beta=\begin{cases}(2u)^{1/(\eta_c+1)} & u\le\frac12\\ \big(\frac{1}{2(1-u)}\big)^{1/(\eta_c+1)} & u>\frac12\end{cases},\qquad u\sim\mathcal U(0,1).
\]
\paragraph{Mutation polynomiale.} Avec probabilité $1/n_{var}$ par variable, un petit saut de loi en puissance, d'amplitude réglée par $\eta_m$.

\begin{figure}[htbp]
\centering
% figure générée par tools/make_figures.py : operateurs
\begin{tikzpicture}
\begin{groupplot}[group style={group size=2 by 1,horizontal sep=1.6cm},width=0.5\linewidth,height=5.2cm,
  title style={font=\small},tick label style={font=\scriptsize},xmin=-0.5,xmax=7.5,xtick={0,...,7},ymin=0]
\nextgroupplot[title={SBX ($\eta_c=15$), parents 2 et 5},xlabel={valeur de l'enfant (indice relâché)},ylabel={densité}]
\addplot[cblue,very thick] coordinates {(-0.45,0) (-0.35,0) (-0.25,0) (-0.15,0) (-0.05,0) (0.05,0) (0.15,0) (0.25,0) (0.35,0.000125) (0.45,0) (0.55,0) (0.65,0.000125) (0.75,0.000125) (0.85,0) (0.95,0.000125) (1.05,0.00025) (1.15,0.000875) (1.25,0.0015) (1.35,0.002625) (1.45,0.007875) (1.55,0.016625) (1.65,0.0375) (1.75,0.10125) (1.85,0.27237) (1.95,0.8075) (2.05,3.3352) (2.15,0.29075) (2.25,0.093) (2.35,0.02425) (2.45,0.005625) (2.55,0.00175) (2.65,0.000375) (2.75,0.000125) (2.85,0) (2.95,0) (3.05,0) (3.15,0) (3.25,0) (3.35,0) (3.45,0) (3.55,0) (3.65,0) (3.75,0) (3.85,0) (3.95,0) (4.05,0) (4.15,0) (4.25,0.000125) (4.35,0.000375) (4.45,0.00175) (4.55,0.005625) (4.65,0.02425) (4.75,0.093) (4.85,0.29075) (4.95,0.83263) (5.05,3.3101) (5.15,0.27238) (5.25,0.10125) (5.35,0.0375) (5.45,0.016625) (5.55,0.007875) (5.65,0.002625) (5.75,0.0015) (5.85,0.000875) (5.95,0.00025) (6.05,0.000125) (6.15,0) (6.25,0.000125) (6.35,0.000125) (6.45,0) (6.55,0) (6.65,0.000125) (6.75,0) (6.85,0) (6.95,0) (7.05,0) (7.15,0) (7.25,0) (7.35,0) (7.45,0)};
\addplot[only marks,mark=triangle*,mark size=3pt,corange] coordinates {(2,0) (5,0)};
\nextgroupplot[title={mutation polynomiale ($\eta_m=20$), parent 3},xlabel={valeur après mutation}]
\addplot[cgreen,very thick] coordinates {(-0.45,0.00025) (-0.35,0) (-0.25,0) (-0.15,0.00025) (-0.05,0) (0.05,0) (0.15,0.0005) (0.25,0) (0.35,0.00075) (0.45,0.001) (0.55,0.0005) (0.65,0.00075) (0.75,0.002) (0.85,0.00225) (0.95,0.00225) (1.05,0.00475) (1.15,0.008) (1.25,0.008) (1.35,0.015) (1.45,0.01825) (1.55,0.0285) (1.65,0.03) (1.75,0.0405) (1.85,0.061) (1.95,0.086) (2.05,0.096) (2.15,0.14325) (2.25,0.19575) (2.35,0.25) (2.45,0.32675) (2.55,0.415) (2.65,0.518) (2.75,0.687) (2.85,0.9075) (2.95,1.1618) (3.05,1.1697) (3.15,0.89025) (3.25,0.716) (3.35,0.544) (3.45,0.388) (3.55,0.31075) (3.65,0.23475) (3.75,0.19) (3.85,0.1425) (3.95,0.10075) (4.05,0.076) (4.15,0.062) (4.25,0.042) (4.35,0.0315) (4.45,0.02425) (4.55,0.022) (4.65,0.012) (4.75,0.009) (4.85,0.0075) (4.95,0.00475) (5.05,0.00225) (5.15,0.00275) (5.25,0.0015) (5.35,0.0015) (5.45,0.00125) (5.55,0.0005) (5.65,0) (5.75,0.00075) (5.85,0) (5.95,0.00025) (6.05,0) (6.15,0) (6.25,0) (6.35,0) (6.45,0) (6.55,0) (6.65,0) (6.75,0) (6.85,0) (6.95,0) (7.05,0) (7.15,0) (7.25,0) (7.35,0) (7.45,0)};
\addplot[only marks,mark=triangle*,mark size=3pt,corange] coordinates {(3,0)};
\end{groupplot}
\end{tikzpicture}

\caption{Distribution des enfants pour les réglages usuels ($\eta_c=15$, $\eta_m=20$). Après arrondi, l'enfant SBX reprend la valeur d'un des parents dans 99,7~\% des cas, et la mutation laisse l'indice inchangé dans 74~\% des cas, sinon le décale de $\pm1$.}\label{fig:operateurs}
\end{figure}

\begin{piege}
Avec des variables à 8 niveaux, l'arrondi transforme ces opérateurs continus. SBX avec $\eta_c=15$ devient un \emph{échange de gènes} (chaque enfant hérite la valeur de l'un ou l'autre parent, variable par variable), et c'est la mutation qui explore, par pas de $\pm1$. Ce n'est pas un défaut, mais il faut le savoir : régler $\eta_c$ ne change presque rien ici.
\end{piege}

\paragraph{Doublons.} Un enfant identique à un individu de la population, ou à un autre enfant, est rejeté avant d'être simulé. Sans cela, la population se remplirait de copies d'un même bon profil, et des simulations seraient gaspillées.

\section{Vérification contre pymoo}
Le script \code{optim/check\_pymoo.py} lance l'implémentation de pymoo (installée dans \code{Code/.venv}) sur le même problème compté, avec les mêmes réglages ($N=16$, SBX $\eta_c=15$, mutation $\eta_m=20$, arrondi, doublons éliminés). Sur 20 graines, le nombre médian de simulations pour retrouver tout le front est :
\begin{center}
\begin{tabular}{lcc}
\toprule
 & notre NSGA-II & pymoo \\
\midrule
S1 (7 profils) & 106,5 & 95,0 \\
S2 (15 profils) & 106,0 & 110,5 \\
\bottomrule
\end{tabular}
\end{center}
Les deux implémentations se comportent de la même façon, aux fluctuations entre graines près.

\exercices
\begin{exercice}
Calculer la distance d'encombrement des trois points d'un front $(0;\,1)$, $(0{,}3;\,0{,}5)$, $(1;\,0)$. Lequel serait éliminé en premier si l'on devait n'en garder que deux ? Pourquoi les extrémités reçoivent-elles $+\infty$ ?
\end{exercice}
\begin{exercice}
Avec $\eta_c=15$ et des parents d'indices 2 et 5, quelle valeur de $\beta$ faudrait-il pour qu'un enfant arrondi soit égal à 3 ? Avec quelle probabilité cela arrive-t-il ?
\end{exercice}

\begin{pistes}
\begin{itemize}
\item Utiliser des opérateurs \emph{entiers} natifs (croisement uniforme, mutation par décalage de $\pm1$ ou par tirage), plus lisibles que SBX arrondi.
\item Utiliser un codage binaire (les 9 bits du QUBO) avec croisement à un point et mutation par bit : c'est la version « génétique » la plus proche de la formulation quantique.
\item Passer à NSGA-III ou SMS-EMOA si l'on ajoute des objectifs.
\end{itemize}
\end{pistes}
